% Section 5: the C++ implementation, what it checks, and what it measured.
% The numbers in this section are the ones Section 6 explains.
\section{Testing with a C++ implementation}
\label{sec:cpp-testing}

The handout leaves implementation optional. It is done here because two claims
in this report are the kind that ought to be measured rather than asserted: that
generation always returns a spanning tree, and that a heuristic guide changes
how much of the maze a search looks at.

\subsection{What was built}

The code is C++20, built with CMake, compiled with
\texttt{-Wall -Wextra -Wpedantic -Werror}, and lives beside this report in
\texttt{implementation/}.

\begin{center}
\begin{tabular}{@{}ll@{}}
  \toprule
  File & Role \\
  \midrule
  \texttt{cell.hpp}            & a cell, and the two axis operations of
                                 Section~\ref{sec:regions} \\
  \texttt{maze.\{hpp,cpp\}}      & the node records: coordinates and neighbour lists \\
  \texttt{maze\_generator.\{hpp,cpp\}} & Algorithm~\ref{alg:generate-maze}, and
                                 the region tree \\
  \texttt{maze\_solver.\{hpp,cpp\}}    & Algorithm~\ref{alg:search}, with the
                                 estimate as tie-break or priority \\
  \texttt{maze\_invariants.\{hpp,cpp\}} & the properties a finished maze must
                                 have, checked from outside \\
  \texttt{maze\_svg.\{hpp,cpp\}}       & drawing, for looking at a result \\
  \bottomrule
\end{tabular}
\end{center}

The invariants are a separate module on purpose: a generator that checked its
own work would only ever confirm its own assumptions.

\subsection{What is checked}

\begin{description}[leftmargin=*,style=nextline,itemsep=4pt]
  \item[Generation returns a spanning tree.] $WH-1$ doors, connected, no cell
    with more than four. Over ten shapes $\times$ twenty seeds, including
    $1\times1$, $1\times40$ and $40\times1$.
  \item[A corridor is laid down whole.] A $1\times40$ area produces a region
    tree of exactly one leaf and $39$ doors, which is
    Section~\ref{sec:corridor-base} observed rather than argued.
  \item[The same seed gives the same maze.] Required for any measurement below
    to be reproducible.
  \item[Every setting returns the same route.] Plain Dijkstra, both estimates as
    tie-breaks, and both as priorities agree on the path, cell for cell ---
    which is Proposition~\ref{prop:spanning-tree} made visible.
  \item[The route is walkable.] Consecutive cells on it are joined by a door.
  \item[Bad input is refused.] A zero dimension, a door between cells that do
    not touch, a door opened twice, a start outside the grid.
\end{description}

\subsection{What it measured}

Two hundred mazes per size, start at $(0,0)$, exit at the far corner, counting
the cells each search settled before the exit came out of the queue.

\begin{center}
\begin{tabular}{@{}lrrrrrrr@{}}
  \toprule
  & & & \multicolumn{3}{c}{estimate as tie-break} & \multicolumn{2}{c}{as priority (A*)} \\
  \cmidrule(lr){4-6}\cmidrule(lr){7-8}
  Size & Cells & Steps & none & Manhattan & Euclidean & Manhattan & Euclidean \\
  \midrule
  $8\times8$     & $64$    & $22.6$   & $55.6$    & $54.8$    & $54.8$    & $\mathbf{41.0}$    & $45.8$ \\
  $16\times16$   & $256$   & $64.4$   & $219.1$   & $217.7$   & $217.7$   & $\mathbf{172.1}$   & $186.9$ \\
  $32\times32$   & $1024$  & $188.0$  & $879.5$   & $877.2$   & $877.2$   & $\mathbf{754.9}$   & $789.7$ \\
  $64\times64$   & $4096$  & $549.4$  & $3587.1$  & $3584.1$  & $3584.1$  & $\mathbf{3278.8}$  & $3359.0$ \\
  $128\times128$ & $16384$ & $1595.9$ & $14128.5$ & $14124.0$ & $14124.0$ & $\mathbf{13303.5}$ & $13488.9$ \\
  \bottomrule
\end{tabular}
\end{center}

Every one of the $1000$ mazes generated for the table was checked to be a
spanning tree. All of them were.

\subsection{Which estimate to use}

The test was meant to settle Manhattan against Euclidean. It settles rather more
than that.

\begin{description}[leftmargin=*,style=nextline,itemsep=5pt]
  \item[As a tie-break, neither.] The saving is $4.5$ cells in $14\,128$, or
    $0.03\%$, and the two estimates are \emph{exactly} equal in every row. That
    is Proposition~\ref{prop:tiebreak-ceiling} arriving as predicted: the
    tie-break can only reorder the last layer, and both estimates put the exit
    first in it.
  \item[As a priority, Manhattan.] It expands fewer cells than Euclidean at
    every size --- $13\,303$ against $13\,489$ at $128\times128$ --- which is
    the dominance of Section~\ref{sec:shortest-path} showing up as a
    measurement.
  \item[But the gain is small either way.] A* with Manhattan still settles
    $81\%$ of the maze, against $86\%$ for plain Dijkstra. The search looks at
    most of the maze whatever it is told.
\end{description}

\begin{keypoint}
The answer is \keyterm{Manhattan, used as a priority} --- and the more useful
finding is how little any of it buys. Section~\ref{sec:analysis} says why that
is a property of mazes and not of this implementation.
\end{keypoint}
